\documentclass[12pt, aspectratio=169]{beamer}
\usetheme[progressbar=frametitle]{metropolis}
\usecolortheme{whale}
\usepackage[utf8]{inputenc}
\usepackage[T1]{fontenc}
\usepackage{lmodern}
\usepackage{amsmath}
\usepackage{amssymb}
\usepackage{setspace}

\title{Type Checking in Lean}
\subtitle{Type Theory and Formal Verification: A Comprehensive Overview}
\author{Academic Presentation}
\institute{Study of Lean 4 Architecture}
\date{March 11, 2026}

\begin{document}

\begin{frame}
  \titlepage
\end{frame}

\begin{frame}{Outline}
  \tableofcontents[hideallsubsections]
\end{frame}

\section{1. The Epistemological Foundation}

\begin{frame}{Empirical Testing vs. Formal Verification}
\begin{itemize}\setlength{\itemsep}{1.2em}
  \item \textbf{Software Testing}: Provides a localized guarantee. Demonstrates the absence of failure under specific, finite conditions.
  \item \textbf{Limitation}: Intrinsically limited by the impossibility of exploring an infinite state space.
  \item \textbf{Formal Verification}: Leverages mathematical proofs to guarantee the absolute impossibility of failure across all defined states.
  \item Correctness becomes a mathematically proven theorem rather than a statistical likelihood.
\end{itemize}
\end{frame}

\begin{frame}{The Lean 4 Vanguard}
\begin{itemize}\setlength{\itemsep}{1.2em}
  \item \textbf{Lean 4}: An interactive theorem prover and functional programming language.
  \item Developed by researchers at Microsoft Research and Carnegie Mellon University.
  \item Fundamentally erases the boundary between a theorem prover and a general-purpose programming language.
  \item Compiles to highly efficient C code safely.
\end{itemize}
\end{frame}

\begin{frame}{Compile-Time Type Inhabitation}
\begin{itemize}\setlength{\itemsep}{1.2em}
  \item In dynamically typed languages (like Python), $2 + 3 = 5$ evaluates to a runtime boolean \texttt{True}.
  \item In Lean, $2 + 3 = 5$ is treated as a complex \textit{compile-time type inhabitation problem}.
  \item Understanding this divergence requires traversing type theory, the Curry-Howard isomorphism, and internal reduction algorithms.
\end{itemize}
\end{frame}

\section{2. A Crash Course in Type Theory}

\begin{frame}{The Untyped $\lambda$-Calculus}
\begin{itemize}\setlength{\itemsep}{1.2em}
  \item \textbf{Untyped $\lambda$-calculus}: Introduced by Alonzo Church (1930s).
  \item Comprises variables ($x$), function abstraction ($\lambda x. t$), and application ($t_1 t_2$).
  \item Turing-complete, but allows unrestricted application.
  \item \textbf{The Problem}: Evaluating a function upon itself can lead to paradoxical, non-terminating terms (mirroring naive set theory inconsistencies).
\end{itemize}
\end{frame}

\begin{frame}{Simply Typed $\lambda$-Calculus ($\lambda^{\to}$)}
\begin{itemize}\setlength{\itemsep}{1.2em}
  \item To introduce discipline and guarantee termination, Church formalized the simply typed lambda calculus.
  \item Every term possesses an intrinsic, immutable type.
  \item Function applications are strictly validated against these types.
  \item If $f : A \to B$ and $x : A$, then $f(x) : B$. Otherwise, it fails immediately.
\end{itemize}
\end{frame}

\begin{frame}{Typing Judgments and Inference Rules}
\begin{itemize}\setlength{\itemsep}{1.2em}
  \item \textbf{Typing Judgment}: $\Gamma \vdash t : A$ 
  \item "Under context $\Gamma$, term $t$ possesses type $A$."
  \item A formal type theory defines mathematical objects via four rules:
  \begin{enumerate}
      \item \textbf{Formation}: How to construct new types.
      \item \textbf{Introduction}: How to construct (inhabit) terms of a new type.
      \item \textbf{Elimination}: How to logically use or consume a term.
      \item \textbf{Computation}: How elimination interacts with introduction (reduction).
  \end{enumerate}
\end{itemize}
\end{frame}

\section{3. Dependent Type Theory (DTT)}

\begin{frame}{The Leap to Dependent Types}
\begin{itemize}\setlength{\itemsep}{1.2em}
  \item Simply typed systems are insufficient for rigorous mathematical reasoning.
  \item \textbf{Dependent Type Theory (DTT)}: Types can dynamically depend on runtime or compile-time values.
  \item \textbf{Example}: A generic \texttt{List} vs. a \texttt{Vector $\alpha$ n}. The type depends on the specific variable $n$ of type \texttt{Nat}, allowing strict array bounds checking at compile time!
\end{itemize}
\end{frame}

\begin{frame}{Paramount Type Constructors in Lean}
\begin{itemize}\setlength{\itemsep}{1.2em}
  \item \textbf{Dependent Product ($\Pi$ type)}:
  \begin{itemize}
      \item Generalizes function types $A \to B$. The return type $B$ depends on input value $x$.
      \item Formally: $\Pi (x : A), B(x)$.
      \item Directly translates to Universal Quantification ($\forall x : A, P(x)$).
  \end{itemize}
  \item \textbf{Dependent Sum ($\Sigma$ type)}:
  \begin{itemize}
      \item Generalizes tuple types $A \times B$. Type of 2nd element depends on 1st element.
      \item Directly translates to Existential Quantification ($\exists x : A, P(x)$).
  \end{itemize}
\end{itemize}
\end{frame}

\begin{frame}{Universes and Paradoxes}
\begin{itemize}\setlength{\itemsep}{1.2em}
  \item If types behave like standard terms, do they have types themselves?
  \item Allowing \texttt{Type : Type} introduces Girard's Paradox (fatal logical inconsistency).
  \item Lean implements a strict, countably infinite hierarchy of universes: \textbf{Calculus of Inductive Constructions (CIC)}.
  \item \texttt{Prop (Sort 0)}: Logical propositions.
  \item \texttt{Type 0 (Sort 1)}: Data types (Nat, Bool).
  \item \texttt{Type 1 (Sort 2)}: Polymorphic types quantifying over \texttt{Type 0}.
\end{itemize}
\end{frame}

\section{4. The Curry-Howard Isomorphism}

\begin{frame}{Propositions as Types}
\begin{itemize}\setlength{\itemsep}{1.2em}
  \item The architectural pillar of interactive theorem proving.
  \item A profound mathematical isomorphism between constructive logic and dependent programming.
  \item \textbf{Mathematical Proposition $P$} $\iff$ \textbf{Type $T$}
  \item \textbf{Formal Proof of $P$} $\iff$ \textbf{Term $t$ inhabiting Type $T$}
  \item Writing a formal proof in Lean is structurally identical to writing a functional program!
\end{itemize}
\end{frame}

\begin{frame}{Logic-to-Type Dictionary}
\begin{table}[]
\begin{tabular}{ll}
\textbf{Constructive Logic}          & \textbf{Dependent Type Theory} \\ \hline
Proposition ($P$)                    & Type ($T$)                     \\
Proof of $P$                         & Term $t$ inhabiting Type $T$   \\
Implication ($P \implies Q$)         & Function Type ($A \to B$)      \\
Universal Quant. ($\forall x, P(x)$) & Dependent Product ($\Pi$ type) \\
Existential Quant. ($\exists x, P(x)$) & Dependent Sum ($\Sigma$ type)  \\
Conjunction ($P \land Q$)            & Product / Tuple Type ($A \times B$) \\
True / Tautology ($\top$)            & Unit Type                      \\
False / Contradiction ($\bot$)       & Empty Type                     \\ \hline
\end{tabular}
\end{table}
\end{frame}

\begin{frame}{Proof Irrelevance}
\begin{itemize}\setlength{\itemsep}{1.2em}
  \item Consider a mathematical proposition $p : Prop$.
  \item If $t_1, t_2 : p$ represent two entirely distinct proof terms asserting the truth of $p$, Lean fundamentally treats $t_1$ and $t_2$ as definitionally equal.
  \item \textbf{Proof Irrelevance} guarantees that proofs carry no excess computational baggage beyond the pure fact that the proposition is true, aiding massive performance scaling.
\end{itemize}
\end{frame}

\section{5. Lean's Pipeline Architecture}

\begin{frame}{Separation of Concerns}
Lean 4 achieves vast scalability by architecturally separating complex heuristic elaboration from the minimal, highly trusted type-checking kernel.
\begin{enumerate}\setlength{\itemsep}{1.0em}
  \item \textbf{Macro Expansion}: Reduces syntactic sugar to core constructs.
  \item \textbf{The Elaborator}: Resource-intensive component that fills in profound gaps, infers implicit arguments, resolves type classes, and generates a fully saturated "pre-term".
\end{enumerate}
\end{frame}

\begin{frame}{The Trusted Kernel}
\begin{itemize}\setlength{\itemsep}{1.2em}
  \item The fully elaborated term is submitted to the trusted C++ kernel.
  \item \textbf{The Kernel is absolutely minimalist}. It represents the final source of truth.
  \item It performs no automation, typeclass inference, or heuristic deduction.
  \item It solely and mechanically verifies that the term complies with the fundamental formation and computation rules of CIC.
\end{itemize}
\end{frame}

\section{6. Machinery of Reduction}

\begin{frame}{Propositional vs. Definitional Equality}
\begin{itemize}\setlength{\itemsep}{1.2em}
  \item \textbf{Propositional Equality}: An explicit mathematical statement (\texttt{Eq a b}). Requires algebraic rewrites via theorems to generate an explicit proof term. (e.g., $x + y = y + x$).
  \item \textbf{Definitional Equality}: A weaker but automated syntactic relationship. Two terms compute to the exact same normal form via Lean's algorithms. (e.g., $2 + 3$ is definitionally equal to $5$).
\end{itemize}
\end{frame}

\begin{frame}{Reduction Algorithms}
To establish definitional equality autonomously, the kernel utilizes:
\begin{itemize}\setlength{\itemsep}{1.2em}
  \item \textbf{$\beta$-Reduction}: Core mechanism of lambda calculus. Substitutes variables in function applications.
  \item \textbf{$\delta$-Reduction}: Mechanically unwinds and expands explicit definitions.
  \item \textbf{$\iota$-Reduction}: Unfolds primitive recursors applied to inductive constructors (powers pattern matching execution).
  \item \textbf{$\zeta$-Reduction}: Expands local \texttt{let} bindings.
\end{itemize}
\end{frame}

\section{7. Python Run-time vs Lean Static Proofs}

\begin{frame}{Dynamic Interpreter Constraints}
\begin{itemize}\setlength{\itemsep}{1.2em}
  \item In Python, \texttt{2 + 3 == 5} is interpreted temporally at execution.
  \item The VM mathematically creates a new object $5$, evaluates $5 == 5$, and pushes a \texttt{True} boolean singleton.
  \item The statement evaluates to a runtime primitive boolean entirely detached from static source code safety.
\end{itemize}
\end{frame}

\begin{frame}{The Propositional Interpretation in Lean}
\begin{itemize}\setlength{\itemsep}{1.2em}
  \item Lean views \texttt{2 + 3 = 5} strictly as a static type declaration: \texttt{Eq Nat (Nat.add 2 3) 5}.
  \item Writing \texttt{rfl} provides the proof term \texttt{Eq.refl \_}.
  \item The Kernel uses $\delta$ and $\iota$ rules to reduce \texttt{Nat.add 2 3} structurally to `Nat.succ (Nat.succ 3)`, which happens to be exactly the integer 5.
  \item Types match, type-check succeeds. The truth is verified completely unconditionally at compile-time.
\end{itemize}
\end{frame}

\begin{frame}{False Type Expressions}
\begin{itemize}\setlength{\itemsep}{1.2em}
  \item In Python, \texttt{2 + 3 == 6} produces the boolean \texttt{False} at runtime to branch logic.
  \item In Lean, writing \texttt{2 + 3 = 6} successfully constructs a perfectly valid type: \texttt{Eq Nat 5 6}\dots
  \item \textbf{But it is inherently uninhabited}. Logical contradiction ($\bot$).
  \item You physically cannot apply a constructor to inhabit it, failing compilation instantly. Truth is forever bound to type inhabitation.
\end{itemize}
\end{frame}

\section{8. Advanced Capabilities}

\begin{frame}{The Equation Compiler \& Metaprogramming}
\begin{itemize}\setlength{\itemsep}{1.2em}
  \item \textbf{Equation Compiler}: Safely translates human-readable recursive pattern matching syntax heavily down into primitive kernel eliminators. It enforces strict structural termination to maintain logic consistency.
  \item \textbf{Lean 4 Metaprogramming}: Exposes internal C++ data structures via Lean itself, allowing users to define perfectly optimized custom syntax and solvers natively without external tools.
\end{itemize}
\end{frame}

\begin{frame}{Tactical Automation \& The Future}
\begin{itemize}\setlength{\itemsep}{1.2em}
  \item Lean heavily leans on advanced tactics like \texttt{simp} via discrimination trees.
  \item Beyond pure algorithms, machine learning integrations like \textbf{Aesop} use neural networks to predict statistically probable proof steps on Mathlib.
  \item The synergy between rigorous formal verification and statistical AI marks the transition into true automated mathematics.
\end{itemize}
\end{frame}

\begin{frame}{Conclusion}
\begin{itemize}\setlength{\itemsep}{1.2em}
  \item Lean is more than a theorem prover; it's a bridge fusing mathematics and software engineering directly onto the Calculus of Inductive Constructions.
  \item Moving past empirical runtime observation into mathematically impenetrable compile-time guarantees is exactly the future of computing.
\end{itemize}
\end{frame}

\begin{frame}
  \begin{center}
      {\Huge \textbf{Thank You}} \\
      \vspace{1em}
      {Questions?}
  \end{center}
\end{frame}

\end{document}
